%\begin{document}

    Chapter 1 sets out the theoretical framework for the thesis. Topics include the background on the Standard Model and kaon physics to provide context for the NA62 experiment and future kaon experiments. 
    The Chapter also include theoretical foundations for the Dark Photon. The NA62 experiment is described in Chapter 2 with descriptions of the sub-detectors and the latest result for the $K^+ \to \pi^+ \nu \bar{\nu}$ decay. 
    During the periods 2021-2024 I completed data taking shifts as well as expert shifts for the KTAG and HLT to assist with the collection of data.

    Chapter 3 describes the development and commissioning of the CEDAR-H detector at CERN. Section 3.1 describes the theoretical background of the CEDAR as well as the basic construction of the vessel, of which is not my own. The design of the CEDAR-H is described in Section 3.1 
    which was completed prior to my involvement. The characterisation of the optical elements was completed by CERN engineers but was consolidated by myself and documented. The simulation data required for the alignment of the spherical mirror was completed by me.
    I was involved in the test-beam in Section 3.3 by completing a number of shifts through out the period. I conducted the analysis of the data from the test-beam shown in this Section. I was also involved in the installation described in Section 3.4 which was conducted at the start of my Long Term Attachment in 2023.
    I was also involved in the commissioning in 2023 described in Section 3.5 and produced all the data analysis shown in this Section. The commissioning in 2024 was also completed by myself and other members of the Birmingham group. I was also involved in the monitoring of the KTAG performance in the periods 2021-2024 by conducing shifts as the detector expert and is briefly described in Section 3.5

    The search of the Dark Photon described in Chapter 4 is entirely my work except for the results referenced from other experiments.

    Chapter 5 set outs the possibility for future kaon experiments. Section 5.1 describes the aims for a High Intensity Kaon Experiments (HIKE) and the planned detector setups for three separate phases. Section 5.2 describes my contributions to the NA62 software were I developed and tested new flexible features to allow the development of new detectors. Section 5.3 follow on with development of an upgraded KTAG which I helped implement into the new software framework. 
    In Section 5.4 the upgraded RICH is described which I completed from design to implementation in the software framework. Section 5.5 provides the detector rates for HIKE Phase 2 which I completed.

    No other qualifications were obtained from the work completed within this thesis.

    %\end{document}